\begin{table}[!htbp]
\centering
\caption{Mexican-Born Immigrants Against All Other Origins, Ages 25 and Over}
\label{tab:mexico_vs_other}
\resizebox{\ifdim\width>\textwidth \textwidth\else \width\fi}{!}{

\begin{tabular}{lcc}
\toprule
 & Mexico & All other origins\\
\midrule
No high school & 52.4 & 14.6\\
High school & 27.7 & 23.5\\
Some college & 17.7 & 44.0\\
Graduate degree & 2.1 & 18.0\\
\midrule
Low-Skill Labor \& Service & 44.0 & 24.4\\
Operatives \& Craft Workers & 30.3 & 12.7\\
Clerical Workers & 7.0 & 10.3\\
Sales Workers & 5.7 & 9.0\\
Managers & 6.6 & 16.2\\
Professional \& Technical Workers & 6.3 & 27.2\\
\midrule
Male & 51.3 & 45.9\\
\bottomrule
\end{tabular}
}
\begin{minipage}{\textwidth}
\footnotesize
\vspace{4pt}
\textit{Notes:} Columns are per cent within origin group, so neither is affected by the sampling design's overrepresentation of the foreign-born. The upper panel is educational attainment, the middle panel occupation and the lower panel sex. Occupation shares are over those who worked in the year observed: \texttt{Occup\_Group} assigns Not in Labor Force on the basis of \texttt{WORKEDYR}, which places roughly a quarter of each origin group outside the occupational distribution, and they are excluded here so the shares describe workers. Military occupations are omitted from the panel; they account for under 0.2 per cent of either group.
\end{minipage}
\end{table}
